\chapter{Đánh giá}
\label{chap:eval}

\section{Bộ dữ liệu đánh giá}

Đánh giá hiện tại dùng \texttt{data/eval/gold\_synthetic\_205.jsonl}. Benchmark gần nhất chạy 100 mẫu trên tập này để kiểm tra chất lượng extraction ở mức baseline. \texttt{data/eval/gold\_smoke.jsonl} vẫn được giữ như smoke test nhỏ cho CI, nhưng không còn là số liệu chính để kết luận baseline.

Model được đánh giá là Qwen2.5-3B base chạy qua Ollama, chưa fine-tune. Vì vậy kết quả bên dưới là baseline trước khi sử dụng dữ liệu feedback thật để cải thiện mô hình hoặc prompt.

\section{Kết quả định lượng}

\begin{table}[H]
\centering
\caption{Kết quả baseline trên benchmark 100 mẫu}
\begin{tabular}{lcc}
\toprule
\textbf{Metric} & \textbf{Kết quả} & \textbf{Gate} \\
\midrule
Precision & 0.8604 & hard: $\geq$ 0.70 \\
Recall & 0.6665 & watch \\
F1 & 0.6886 & no drop $>$ 0.05 \\
Assignee accuracy & 0.5232 & watch \\
Hallucination rate & 0.0\% & hard: no $>$2pp regression \\
Schema failure rate & 0.0\% & hard: no regression \\
\bottomrule
\end{tabular}
\end{table}

Kết quả cho thấy precision, schema validity và hallucination control đã đủ tốt cho baseline hiện tại. Recall và assignee accuracy vẫn là watch metrics cần cải thiện bằng dữ liệu feedback, prompt tuning hoặc fine-tuning. Promotion policy hiện không yêu cầu mọi metric vượt một ngưỡng cao tuyệt đối; candidate được so tương đối với baseline để tránh regression.

\begin{figure}[H]
\centering
\begin{tikzpicture}
\begin{axis}[
  ybar,
  bar width=18pt,
  width=0.9\textwidth,
  height=5.2cm,
  ymin=0,
  ymax=1.05,
  ylabel={Score},
  symbolic x coords={Precision,Recall,F1},
  xtick=data,
  nodes near coords,
  nodes near coords align={vertical},
  every node near coord/.append style={font=\small},
  ymajorgrids=true,
  grid style={lightgray}
]
\addplot+[fill=primary] coordinates {(Precision,0.8604) (Recall,0.6665) (F1,0.6886)};
\end{axis}
\end{tikzpicture}
\caption{Kết quả baseline của Qwen2.5-3B trên benchmark 100 mẫu}
\label{fig:baseline-metrics}
\end{figure}

\section{Cách tính metric}

Evaluation harness đo chất lượng extraction ở mức action item. Mỗi sample gồm transcript turns, roster và danh sách action items gold. Hệ thống chạy lại hàm extraction, chuyển kết quả Pydantic sang JSON, sau đó so sánh từng predicted task với gold task.

Matching mô tả task dùng BM25 trên trường \texttt{description}; nếu điểm BM25 không vượt ngưỡng, code fallback sang Jaccard overlap để tránh bỏ sót các câu ngắn có từ vựng trùng rõ ràng. Với mỗi predicted task, nếu match với một gold task chưa được dùng, hệ thống tính là true positive; predicted task không match là false positive; gold task không được match là false negative. Precision, recall và F1 được tính từ các giá trị này. Assignee accuracy chỉ được tính trên các true positive có cả assignee dự đoán và assignee gold.

Schema failure rate đo số sample mà pipeline không trả được output hợp lệ, ví dụ lỗi parse JSON hoặc lỗi guardrail nghiêm trọng làm extraction thất bại. Hallucination trong evaluation được ước lượng bằng overlap giữa từ nội dung trong task description và transcript; nếu overlap thấp hơn ngưỡng, predicted task bị xem là không có bằng chứng đủ mạnh trong transcript. Các metric này chưa thay thế đánh giá thủ công, nhưng đủ để phát hiện regression khi thay prompt, model hoặc dữ liệu training.

\section{Kiểm thử vận hành}

Các kiểm tra đã được dùng để xác nhận hành vi hệ thống:

\begin{itemize}[leftmargin=*]
  \item Ruff check cho mã Python.
  \item Pytest cho các phần STT, retrain promotion, deploy promoted model và export.
  \item Next.js lint/build cho frontend.
  \item Docker Compose config cho profile mặc định và profile \texttt{mlops}.
  \item Health check API/web sau khi build lại container.
\end{itemize}

\section{Đánh giá định tính}

UI Human Feedback sau chỉnh sửa giúp người dùng hiểu ngữ cảnh multi-speaker rõ hơn: transcript và speaker evidence được đặt cạnh speaker mapping, task evidence hiển thị đoạn hội thoại nguồn. Điều này giảm rủi ro gán nhầm người phụ trách khi audio có nhiều speaker.

Grafana/Prometheus phản ánh trạng thái worker, meeting pipeline và số liệu feedback. Các panel hiện tại đủ để quan sát pipeline trong phạm vi dự án; những chỉ số vận hành nâng cao như latency percentile, queue age hoặc GPU utilization có thể bổ sung khi mở rộng tải hoặc bổ sung GPU training.
